% Figure from MATH 590 notes, section 29. Compile with the notes preamble.
\begin{tikzpicture}[scale=1.15,
  mid/.style={postaction={decorate}, decoration={markings, mark=at position 0.58 with {\arrow{Stealth[length=2.6mm]}}}}]
  \fill[black!4] (-0.612,-1.478) -- (0.612,-1.478) -- (1.478,-0.612) -- (1.478,0.612) -- (0.612,1.478) -- (-0.612,1.478) -- (-1.478,0.612) -- (-1.478,-0.612) -- cycle;
  \draw[black!45, dashed, thin] (-0.612,-1.478) -- (0.612,1.478);
  \draw[figB, thick, mid] (-0.612,-1.478) -- (0.612,-1.478);
  \node[figB, font=\small] at (0.000,-1.720) {$a_1$};
  \draw[figR, thick, mid] (0.612,-1.478) -- (1.478,-0.612);
  \node[figR, font=\small] at (1.216,-1.216) {$b_1$};
  \draw[figB, thick, mid] (1.478,0.612) -- (1.478,-0.612);
  \node[figB, font=\small] at (1.720,0.000) {$a_1$};
  \draw[figR, thick, mid] (0.612,1.478) -- (1.478,0.612);
  \node[figR, font=\small] at (1.216,1.216) {$b_1$};
  \draw[figB, thick, mid] (0.612,1.478) -- (-0.612,1.478);
  \node[figB, font=\small] at (0.000,1.720) {$a_2$};
  \draw[figR, thick, mid] (-0.612,1.478) -- (-1.478,0.612);
  \node[figR, font=\small] at (-1.216,1.216) {$b_2$};
  \draw[figB, thick, mid] (-1.478,-0.612) -- (-1.478,0.612);
  \node[figB, font=\small] at (-1.720,0.000) {$a_2$};
  \draw[figR, thick, mid] (-0.612,-1.478) -- (-1.478,-0.612);
  \node[figR, font=\small] at (-1.216,-1.216) {$b_2$};
  \fill[black] (-0.612,-1.478) circle (1.6pt);
  \fill[black] (0.612,-1.478) circle (1.6pt);
  \fill[black] (1.478,-0.612) circle (1.6pt);
  \fill[black] (1.478,0.612) circle (1.6pt);
  \fill[black] (0.612,1.478) circle (1.6pt);
  \fill[black] (-0.612,1.478) circle (1.6pt);
  \fill[black] (-1.478,0.612) circle (1.6pt);
  \fill[black] (-1.478,-0.612) circle (1.6pt);
  \node[font=\scriptsize, black!60] at (0.6,-0.55) {$a_1b_1a_1^{-1}b_1^{-1}$};
  \node[font=\scriptsize, black!60] at (-0.6,0.55) {$a_2b_2a_2^{-1}b_2^{-1}$};
\end{tikzpicture}
